\documentclass[a4paper,11pt]{article}

\usepackage{amsmath}
\usepackage{url}
\usepackage{hyperref}
\usepackage{enumitem}

\title{Milestone 3: Critical Evaluation and Transfer}
\author{Christian Percival}
\date{}

\begin{document}

\maketitle

\section*{Purpose}

\begin{itemize}
    \item This milestone evaluates the strengths, limitations, risks, assumptions, and applicability of overrun handling.
    \item The goal is not to say whether the paper is simply good or bad.
    \item The goal is to understand where the approach works, where it may fail, and how it transfers to realistic embedded systems.
\end{itemize}

\section*{Strengths of the Approach}

\begin{itemize}
    \item Handles variable execution times more efficiently than pure WCET scheduling.
    \item Allows higher normal-case task frequencies.
    \item Improves control performance when normal execution time is much smaller than WCET.
    \item Uses bandwidth reservation to limit interference between tasks.
    \item Provides temporal isolation.
    \item Maintains hard real-time guarantees under bounded worst-case assumptions.
    \item CBShd improves standard CBS by handling partial overruns more accurately.
    \item The approach is strongly motivated by realistic control applications such as visual tracking.
\end{itemize}

\section*{Limitations}

\begin{itemize}
    \item Requires a known or bounded WCET.
    \item WCET estimation is difficult on modern processors.
    \item Cache, interrupts, operating system behavior, and background processes can affect timing.
    \item Requires knowledge of normal computation time.
    \item Normal computation time may change depending on input data.
    \item Remaining computation time may be difficult to estimate accurately.
    \item Server-based scheduling introduces runtime overhead.
    \item Implementation requires real-time scheduling support.
    \item A simple desktop operating system cannot provide the same guarantees as a real-time kernel.
\end{itemize}

\section*{Risks and Hidden Assumptions}

\begin{itemize}
    \item If WCET is underestimated, hard deadlines may still be missed.
    \item If normal execution time is underestimated, overruns may occur too often.
    \item If server budget is too small, overhead may become too high.
    \item If server budget is too large, overrun handling becomes less precise.
    \item If a task overruns while holding a shared resource, blocking problems can occur.
    \item The theory assumes bounded behavior, but real image processing can have input-dependent complexity.
\end{itemize}

\section*{Runtime Implications}

\begin{itemize}
    \item Budget management costs CPU time.
    \item Deadline postponement costs CPU time.
    \item Smaller chunks lead to more frequent budget exhaustion.
    \item More frequent budget exhaustion leads to more scheduler overhead.
    \item The paper reports that overhead must be considered when choosing server budget.
\end{itemize}

\section*{Memory and Scalability Implications}

\begin{itemize}
    \item Each task needs server state:
    \begin{itemize}
        \item current budget,
        \item server period,
        \item current deadline,
        \item remaining computation estimate.
    \end{itemize}
    \item Memory overhead is probably small for a small number of tasks.
    \item Analysis becomes more complex as the number of tasks increases.
    \item Resource sharing makes the analysis more complex.
    \item More tasks mean more possible interactions and more scheduling decisions.
\end{itemize}

\section*{Suitable Application Examples}

\begin{itemize}
    \item Robotic visual tracking:
    \begin{itemize}
        \item normal case: object found quickly,
        \item overrun case: larger search region needed.
    \end{itemize}
    \item Radar tracking:
    \begin{itemize}
        \item normal case: target found near predicted position,
        \item overrun case: wider search required.
    \end{itemize}
    \item Sensor fusion:
    \begin{itemize}
        \item normal case: simple sensor update,
        \item overrun case: additional filtering or correction required.
    \end{itemize}
    \item Autonomous systems:
    \begin{itemize}
        \item normal case: predictable environment,
        \item overrun case: unexpected obstacle or uncertain sensor data.
    \end{itemize}
\end{itemize}

\section*{Unsuitable or Problematic Application Examples}

\begin{itemize}
    \item Systems where WCET cannot be bounded realistically.
    \item Systems where tasks must not be delayed after an overrun.
    \item Systems with extremely limited processor power where server overhead is too expensive.
    \item Systems where object detection uses complex AI models with highly unpredictable execution time.
    \item Systems where shared resources are held for long unpredictable times.
    \item Very simple microcontroller applications where fixed WCET scheduling is already sufficient.
\end{itemize}

\section*{Transfer to Student Prototype}

\begin{itemize}
    \item A full real-time kernel implementation is probably too large for this seminar.
    \item A controlled simulation is more realistic.
    \item The prototype can demonstrate:
    \begin{itemize}
        \item normal execution time,
        \item occasional overruns,
        \item WCET-based inefficiency,
        \item deadline misses without overrun handling,
        \item improvement using simplified budget/deadline handling.
    \end{itemize}
    \item Python and OpenCV can be used for the demonstration.
    \item Synthetic images can be generated to control easy, medium, hard, and worst-case inputs.
\end{itemize}

\section*{Possible Extensions}

\begin{itemize}
    \item Implement a simple image-processing overrun demo.
    \item Compare WCET-based scheduling, normal-budget scheduling, and simplified overrun handling.
    \item Add plots:
    \begin{itemize}
        \item execution time per frame,
        \item deadline misses per strategy,
        \item CPU utilization per strategy,
        \item average frequency per strategy.
    \end{itemize}
    \item Later extension:
    \begin{itemize}
        \item run the same experiment on Raspberry Pi,
        \item compare desktop timing and embedded Linux timing,
        \item measure timing jitter.
    \end{itemize}
\end{itemize}

\section*{Critical Conclusion}

\begin{itemize}
    \item The approach is strong because it addresses a real weakness of WCET-only scheduling.
    \item It is especially useful when normal execution time is much smaller than WCET.
    \item It is less useful when WCET is close to normal execution time.
    \item The main practical challenge is accurate timing estimation and implementation overhead.
    \item For the seminar, a simplified prototype is suitable as a demonstration, but it should not be presented as a full real-time scheduler.
\end{itemize}

\end{document}